% !TEX program = lualatex
\documentclass[11pt, letterpaper]{article}

\usepackage{modern}
\addbibresource{references.bib}

\begin{document}

\section{First section}
This document uses the Palatino-like TeX Gyre Pagella font for the main text and
math, and the Optima font for headings. Optima ships as a system font on macOS
but is proprietary and is not bundled with TeX Live (including
\texttt{texlive-full} on Debian); on Linux, substitute URW Classico (CTAN
package \texttt{classico}), which is a metric-compatible free clone. The line
spacing is set to 1.2, and the margins are configured to be slightly less wide
than the default. All of this lives in \texttt{modern.sty} --- this file just
loads the style with \verb|\usepackage{modern}|.

$$ \int_0^\infty e^{-x^2}\, dx = \frac{\sqrt{\pi}}{2} $$

Aa Bb Cc Dd Ee Ff Gg Hh Ii Jj Kk Ll Mm Nn Oo Pp Qq Rr Ss Tt Uu Vv Ww Xx Yy Zz
1234567890

\emph{This is emphasized text.}

\section{Compiling}
This template is built with Lua\LaTeX{} (required for \texttt{fontspec} and
\texttt{unicode-math}) and uses \texttt{biblatex} with the Biber backend for
the bibliography. From the directory containing \texttt{modern.tex},
\texttt{modern.sty}, and \texttt{references.bib}, run:

\begin{verbatim}
lualatex modern.tex
biber modern
lualatex modern.tex
lualatex modern.tex
\end{verbatim}

The first \texttt{lualatex} pass writes the citation data, \texttt{biber}
compiles the bibliography, and the final two \texttt{lualatex} passes resolve
cross-references. The equivalent one-liner with \texttt{latexmk} is:

\begin{verbatim}
latexmk -lualatex modern.tex
\end{verbatim}

\section{Body example}
Modern empirical industrial organization emphasizes the structural modeling of
markets \citep{ackerberg2007econometric}. The aim is to closely represent the
primitive factors and the strategic conduct of agents---on both demand and
production sides---that ultimately determine market outcomes. The structural
approach is often contrasted with reduced-form methods, which are more agnostic
about the underlying economic model and focus on estimating relationships
between observed variables \citep{angrist2010credibility}. Structural models can
be used to conduct counterfactual analysis and to predict the effects of policy
changes, while reduced-form models are typically more focused on identifying
causal relationships in the data. \citet{nevo2001measuring} provides a canonical
example of the structural approach in the ready-to-eat cereal industry.

\printbibliography

\end{document}
